% ============================================================================
% CAPÍTULO 16: EL KERNEL C++20 — ANATOMÍA DEL MOTOR DE SILICIO
% ============================================================================
\chapter{El Kernel C++20: Anatomía del Motor de Silicio}


\textit{El código no miente. Los comentarios mienten.\\
Los benchmarks sin código fuente adjunto mienten más.} --- --- Regla 10, Constitución POLYDIM

\section{Filosofía del Kernel: Correctness Before Performance}

El kernel C++20 de POLYDIM (\texttt{kernel\_cpp\_v762.cpp}, 552 líneas) encarna
una filosofía que invierte el orden habitual de prioridades en computación de alto
rendimiento: \textit{corrección aritmética primero, luego optimización}.

Esta filosofía emerge directamente de los 300+ ciclos de Red Team que descubrieron
bugs críticos precisamente en el código ``optimizado'' que sacrificaba validaciones
en favor de rendimiento. Los ejemplos paradigmáticos:

\begin{itemize}
\item El umbral mágico \texttt{1e-15} que rechazaba vectores perfectamente válidos
de norma \texttt{1e-200} (Ciclo 38).
\item El cálculo de norma con \texttt{np.linalg.norm} que desbordaba a \texttt{inf}
para vectores de norma \texttt{1e305} (Ciclo 29 --- el oráculo del test más débil
que el código bajo test).
\item El \texttt{static\_cast<double>} sin comprobación que convertía silenciosamente
subnormales en exactamente 0.0 (Ciclo 22).
\end{itemize}

\section{La Taxonomía de Errores Numéricos IEEE-754}

Antes de analizar el código, es necesario formalizar la taxonomía de errores que
el kernel debe detectar y manejar:

\begin{definition}[Clases de Números IEEE-754 Double Precision]
El estándar IEEE-754 define las siguientes clases de números de doble precisión:
\begin{align}
\text{Normal:} &\quad 2^{-1022} \le |x| \le (2 - 2^{-52}) \cdot 2^{1023} \\
\text{Subnormal:} &\quad 0 < |x| < 2^{-1022} \\
\text{Zero:} &\quad x = \pm 0 \\
\text{Infinito:} &\quad x = \pm\infty \\
\text{NaN:} &\quad x \text{ (Not a Number)}
\end{align}
\end{definition}

\begin{theorem}[Propagación de NaN]
Si $x = \text{NaN}$, entonces para toda operación aritmética $\circ$:
$x \circ y = \text{NaN}$ para todo $y$. En particular, \texttt{NaN < NaN} es
\texttt{false}, \texttt{NaN == NaN} es \texttt{false} (¡solo caso donde $x \ne x$!).
\end{theorem}

La prueba del Red Team exploitó exactamente esta propiedad: un NaN silencioso
en el vector de entrada propagaría a todos los cálculos de norma produciendo
una norma \texttt{NaN}, que luego comparada con un umbral real produciría
\texttt{false} (``NaN $\le$ threshold''), haciendo pasar la validación cuando
debería fallar.

\section{El Acumulador de Neumaier: Suma Compensada en Alta Dimensión}

El problema de sumar $D$ números de punto flotante con $D = 10^6$ es
numéricamente no trivial. La suma naïve acumula un error $\mathcal{O}(D\varepsilon_{\text{mach}})$,
donde $\varepsilon_{\text{mach}} = 2^{-52} \approx 2.22 \times 10^{-16}$.

Para $D = 10^6$:
\begin{equation}
\text{Error naïve} \approx D \cdot \varepsilon_{\text{mach}} = 10^6 \times 2.22 \times 10^{-16} = 2.22 \times 10^{-10}
\end{equation}

Este error es inadmisible para la verificación de ortonormalidad donde requerimos
$\|{Y^\top Y - I}\|_{\max} \le 10^{-12}$.

\begin{algorithm}


\begin{algorithmic}[1]
\Procedure{NeumaierSum}{$a_1, a_2, \ldots, a_D$}
  \State $s \gets 0.0$; $c \gets 0.0$ \Comment{$s$: suma principal, $c$: compensación}
  \For{$i = 1$ to $D$}
    \State $t \gets s + a_i$
    \If{$|s| \ge |a_i|$}
      \State $c \gets c + (s - t) + a_i$ \Comment{$a_i$ pequeño: $(s - t) + a_i$ exacto}
    \Else
      \State $c \gets c + (a_i - t) + s$ \Comment{$s$ pequeño: $(a_i - t) + s$ exacto}
    \EndIf
    \State $s \gets t$
  \EndFor
  \State \Return $s + c$
\EndProcedure
\end{algorithmic}
\end{algorithm}

\begin{theorem}[Error del Acumulador de Neumaier]

El error del algoritmo de Neumaier satisface:
\begin{equation}
\left|\text{NeumaierSum}(a_1, \ldots, a_D) - \sum_{i=1}^D a_i\right| \le (2\varepsilon_{\text{mach}} + \mathcal{O}(\varepsilon_{\text{mach}}^2)) \max_i |a_i|
\end{equation}
Es decir, el error es \textbf{independiente de $D$}, en contraste con el error $\mathcal{O}(D\varepsilon_{\text{mach}})$ de la suma naïve.
\end{theorem}

La implementación en el kernel V762:

\begin{verbatim}
struct alignas(64) NeumaierAcc {
    double sum;
    double c;

    inline void init() { sum = 0.0; c = 0.0; }

    inline void add(double val) {
        double t = sum + val;
        if (std::abs(sum) >= std::abs(val)) {
            c += (sum - t) + val;  // val pequeño: round-off exacto
        } else {
            c += (val - t) + sum;  // sum pequeño: round-off exacto
        }
        sum = t;
    }

    inline double total() const { return sum + c; }
};
\end{verbatim}

El \texttt{alignas(64)} garantiza que el acumulador ocupa una única línea de caché
(64 bytes), eliminando el \textit{false sharing} entre acumuladores de distintos hilos OpenMP.

\section{El Algoritmo de Rodrigues de 2 Pasadas}

La función principal del kernel --- \texttt{polydim\_apply\_rodrigues\_geodesic\_f64}
--- implementa la rotación geodésica sobre $S^{D-1}$ mediante un algoritmo de
2 pasadas sobre el vector $\mathbf{y} \in \mathbb{R}^D$:

\begin{align}
\text{Pasada 1 (productos punto compensados):} &\quad
\langle\mathbf{y}, \mathbf{u}\rangle, \langle\mathbf{y}, \mathbf{v}\rangle,
\langle\mathbf{u}, \mathbf{u}\rangle, \langle\mathbf{v}, \mathbf{v}\rangle,
\langle\mathbf{u}, \mathbf{v}\rangle \\
\text{Pasada 2 (actualización FMA):} &\quad
y_{\text{out},i} = \text{fma}(\alpha, u_i, \text{fma}(\beta, v_i, y_i))
\end{align}

donde:
\begin{align}
\alpha &= -\operatorname{vers}(\theta)\langle\mathbf{y}, \mathbf{u}\rangle - \sin(\theta)\langle\mathbf{y}, \mathbf{v}\rangle \\
\beta  &= -\operatorname{vers}(\theta)\langle\mathbf{y}, \mathbf{v}\rangle + \sin(\theta)\langle\mathbf{y}, \mathbf{u}\rangle \\
\operatorname{vers}(\theta) &= 2\sin^2(\theta/2) = 1 - \cos(\theta) \quad \text{(estable para } \theta \to 0\text{)}
\end{align}

\begin{verbatim}
POLYDIM_EXPORT int32_t POLYDIM_CALL 
polydim_apply_rodrigues_geodesic_f64(
    const double* y,
    const double* __restrict u,
    const double* __restrict v,
    double* y_out,
    double theta,
    uint64_t D
) {
    if (!y || !u || !v || !y_out) return POLYDIM_ERR_NULL_POINTER;
    if (D == 0) return POLYDIM_ERR_INVALID_DIMENSION;

    FtzDazGuard guard;  // Proteger registros FPU

    int max_threads = omp_get_max_threads();
    std::vector<NeumaierAcc> acc_yu(max_threads);
    std::vector<NeumaierAcc> acc_yv(max_threads);
    std::vector<NeumaierAcc> acc_uu(max_threads);
    std::vector<NeumaierAcc> acc_vv(max_threads);
    std::vector<NeumaierAcc> acc_uv(max_threads);
    // Inicializar todos los acumuladores
    for (int t = 0; t < max_threads; ++t) {
        acc_yu[t].init(); acc_yv[t].init();
        acc_uu[t].init(); acc_vv[t].init(); acc_uv[t].init();
    }

    int nan_detected = 0;
    int subnormal_detected = 0;

    // PASADA 1: Suma compensada de productos punto + detección de patologías
    #pragma omp parallel reduction(|:nan_detected,subnormal_detected)
    {
        FtzDazGuard thread_guard;  // Por hilo: no contaminar otros hilos
        int tid = omp_get_thread_num();
        NeumaierAcc local_yu; local_yu.init();
        // [... inicialización similar para los demás ...]

        #pragma omp for schedule(static)
        for (int64_t i = 0; i < static_cast<int64_t>(D); ++i) {
            double yi = y[i]; double ui = u[i]; double vi = v[i];
            // Detección de NaN/Inf (propagación explícita)
            if (std::isnan(yi)||std::isnan(ui)||std::isnan(vi)||
                std::isinf(yi)||std::isinf(ui)||std::isinf(vi)) {
                nan_detected = 1;
            }
            // Detección de subnormales (SOLO en compiladores que la soporten)
            #if defined(__GNUC__) || defined(__clang__)
            if (std::fpclassify(yi)==FP_SUBNORMAL ||
                std::fpclassify(ui)==FP_SUBNORMAL ||
                std::fpclassify(vi)==FP_SUBNORMAL) {
                subnormal_detected = 1;
            }
            #endif
            local_yu.add(yi * ui);
            // [... local_yv, local_uu, local_vv, local_uv ...]
        }
        acc_yu[tid] = local_yu;
        // [... asignar los demás ...]
    }

    if (nan_detected)       return POLYDIM_ERR_NAN_OR_INF;
    if (subnormal_detected) return POLYDIM_ERR_SUBNORMAL_DETECTED;

    // Reducción final: combinar acumuladores de todos los hilos
    NeumaierAcc total_yu; total_yu.init();
    for (int t = 0; t < max_threads; ++t) {
        total_yu.add(acc_yu[t].total());
        // [... total_yv, total_uu, total_vv, total_uv ...]
    }

    double yu = total_yu.total();
    double uu = total_uu.total();  // ||u||^2
    double vv = total_vv.total();  // ||v||^2

    // VALIDACIÓN: umbral matemáticamente correcto (Ciclo 38)
    // Rechazar SOLO si la norma es denormal (inverso desborda)
    if (uu <= std::numeric_limits<double>::min() ||
        vv <= std::numeric_limits<double>::min()) {
        return POLYDIM_ERR_DEGENERATE_NORM;
    }

    // Versine: 1 - cos(theta) = 2*sin^2(theta/2)
    // Estabilidad numérica para theta -> 0: evita cancelación catastrófica
    double half_theta = 0.5 * theta;
    double sn_half = std::sin(half_theta);
    double vers = 2.0 * sn_half * sn_half;  // *** FIX P1-5: signo correcto ***
    double sn = std::sin(theta);

    // Coeficientes de actualización
    double alpha = -vers * yu - sn * yv;   // FIX #1: orientación exacta
    double beta  = -vers * yv + sn * yu;   // FIX #1: orientación exacta

    // PASADA 2: Actualización streaming con std::fma (2 redondeos en vez de 4)
    // y_out[i] = y[i] + alpha*u[i] + beta*v[i]
    // = fma(alpha, u[i], fma(beta, v[i], y[i]))
    #pragma omp parallel for schedule(static)
    for (int64_t i = 0; i < static_cast<int64_t>(D); ++i) {
        y_out[i] = std::fma(alpha, u[i], std::fma(beta, v[i], y[i]));
    }

    return POLYDIM_SUCCESS;
}
\end{verbatim}

\section{¿Por Qué \texttt{std::fma} Reduce el Error a la Mitad?}

La operación \texttt{fma(a, b, c)} computa $a \cdot b + c$ con un único redondeo
final (en vez de dos redondeos separados: uno para $a \cdot b$ y otro para la suma).

Sin FMA:
\begin{align}
\text{round}(a \cdot b) &= a \cdot b \cdot (1 + \delta_1), \quad |\delta_1| \le \varepsilon_{\text{mach}} \\
\text{round}(\text{round}(a \cdot b) + c) &= (a \cdot b \cdot (1 + \delta_1) + c)(1 + \delta_2) \\
\text{Error total} &\le |a \cdot b| \varepsilon_{\text{mach}} + |a \cdot b \cdot (1 + \delta_1) + c| \varepsilon_{\text{mach}}
\end{align}

Con FMA:
\begin{equation}
\text{fma}(a, b, c) = (a \cdot b + c)(1 + \delta), \quad |\delta| \le \varepsilon_{\text{mach}}
\end{equation}

La reducción del número de redondeos de 2 a 1 se traduce en una reducción del error
de aproximadamente un factor 2 para cada operación de la Pasada 2.

\section{Detección de Subnormales: La Bomba de Tiempo}

Los números subnormales son la fuente de errores más insidiosa en computación numérica
de alto rendimiento. El problema no es que los subnormales sean incorrectos: el estándar
IEEE-754 los maneja correctamente. El problema es que:

\begin{enumerate}
\item En hardware x86 con FTZ habilitado (común en HPC), los subnormales se redondean
a 0 silenciosamente, cambiando el resultado matemático.
\item En hardware sin FTZ, las operaciones con subnormales pueden ser 100--1000x más
lentas que con normales (``denormal penalty'').
\item La PMTP Reporte Nocturno (Iteración 2, 01:00 AM) descubrió que multiplicaciones
sucesivas de matrices de reflexión de Clifford en FP16 pueden inducir Underflow
(subnormales) tras millones de rebotes WAN.
\end{enumerate}

La solución de POLYDIM es:
\begin{itemize}
\item Detectar subnormales en la entrada (Pasada 1) y retornar \texttt{POLYDIM\_ERR\_SUBNORMAL\_DETECTED = -4}.
\item Habilitar FTZ/DAZ durante el cálculo interno vía \texttt{FtzDazGuard} para garantizar
rendimiento predecible.
\item Mantener todos los resultados intermedios en FP32 mínimo antes de cuantizar a FP16
(solución del Reporte Nocturno Iteración 2).
\end{itemize}

\section{El Código de Retorno ABI: Contrato Inter-Lenguaje}

Un aspecto crítico del diseño es el ABI (\textit{Application Binary Interface}) de
los códigos de retorno. El Red Team Ciclo 23 documentó que sin un contrato explícito,
cada integrador reinventa (incorrectamente) la semántica de retry:

\begin{verbatim}
enum PolydimStatusCode : int32_t {
    POLYDIM_SUCCESS             =   0,  // Éxito
    POLYDIM_ERR_NULL_POINTER    =  -1,  // Argumento NULL
    POLYDIM_ERR_INVALID_DIM     =  -2,  // D=0 o D > DIM_MAX
    POLYDIM_ERR_NAN_OR_INF      =  -3,  // NaN o Inf en entrada
    POLYDIM_ERR_SUBNORMAL       =  -4,  // Subnormal detectado
    POLYDIM_ERR_INSTABILITY     =  -5,  // Deriva > tolerancia Higham
    POLYDIM_ERR_TOPO_FRAGMENT   =  -6,  // Betti-0 > 1 (enjambre fragmentado)
    POLYDIM_ERR_BUFFER_OVF      =  -7,  // n > 4096 en Betti guard
    POLYDIM_ERR_DEGENERATE_NORM =  -8,  // Norma denormal (inverso desborda)
    POLYDIM_ERR_SEQLOCK_RACE    =  -9,  // SEQLock: secuencia cambió durante lectura
    POLYDIM_ERR_PANIC_CAUGHT    = -99,  // Panic de Rust capturado
    // Códigos Rust-específicos:
    // -12: magic number inválido (nodo corrompido)
    // -13: double-free (free_lock ya tomado)
};
\end{verbatim}

\section{Benchmarks del Kernel C++: Datos Empíricos Certificados}

\begin{tabular}{llll}
\hline
\textbf{Función} & \textbf{D} & \textbf{Tiempo} & \textbf{Drift/Error} & \textbf{Estado} \\
\hline
\texttt{rodrigues\_geodesic} & $10^6$ & $35.06$ ms & $4.44 \times 10^{-16}$ & \textcolor{polydimgreen}{PASS} \\
\texttt{rodrigues\_geodesic} & $10^7$ & $350.6$ ms & $4.44 \times 10^{-16}$ & \textcolor{polydimgreen}{PASS} \\
\texttt{pmtp\_round\_trip} & $10^6$ & $10.69$ ms & $0.000 \times 10^{+0}$ & \textcolor{polydimgreen}{PASS} \\
\texttt{cayley\_smw\_k8} & $10^6$ & $3928.74$ ms & $3.33 \times 10^{-14}$ & \textcolor{polydimgreen}{PASS} \\
\texttt{betti1\_guard} & $n=8$ & $< 1$ ms & N/A & \textcolor{polydimgreen}{PASS} \\
\texttt{slerp\_nan\_input} & $10^6$ & $< 1$ ms & N/A & \textcolor{polydimgreen}{PASS (-3)} \\
\texttt{slerp\_inf\_input} & $10^6$ & $< 1$ ms & N/A & \textcolor{polydimgreen}{PASS (-3)} \\
\texttt{slerp\_zero\_vector} & $10^6$ & $< 1$ ms & N/A & \textcolor{polydimgreen}{PASS (-8)} \\
\texttt{slerp\_subnormal} & $10^6$ & $< 1$ ms & N/A & \textcolor{polydimgreen}{PASS (-4)} \\
\hline


\end{tabular}

\textbf{[CERTIFICADO]} 
\textbf{Certificación V762 --- Exit Code 0:}
Todos los benchmarks fueron ejecutados en el hardware real del investigador
(AMD x64, Windows 11, Python 3.11, NumPy 1.26) el 2026-09-19.
Logs completos disponibles en Apéndice B.
La simulación o generación artificial de datos está estrictamente prohibida
(Regla 10, Constitución POLYDIM).


\section{Arquitectura V772: BLAS Loader Runtime Dinámico y Threadpool Guard}


Para escalar a dimensiones masivas ($D \ge 50{,}000, K = 512$), el kernel V772 implementa un despachador dinámico de álgebra lineal de Nivel 3 (\texttt{cblas\_dsyrk}, \texttt{cblas\_dgemm}) libre de dependencias estáticas de compilación.

\subsection{Carga Dinámica Segura y Smoke Test en Frío}
El cargador (\texttt{LoadLibraryExW} en Windows / \texttt{dlopen} en Linux) busca dinámicamente librerías optimizadas del proveedor (\texttt{mkl\_rt.dll}, \texttt{libopenblas.dll}). Antes de validar los punteros a funciones:
\begin{enumerate}
    \item Ejecuta un producto matricial simétrico $2 \times 2$ de prueba en frío.
    \item Verifica bit a bit que el resultado coincida con la solución exacta IEEE-754.
    \item \textbf{Fallback Transparente:} Si no se detecta ninguna DLL en el sistema, conmuta sin abortar a un micro-kernel interno teselado en bloques de $32 \times 32 \times 32$ con paralelismo OpenMP.
\end{enumerate}

\subsection{Gobernanza de Hilos y Prevención de Oversubscription}
Para erradicar la contención por cambio de contexto en regiones paralelas anidadas, el kernel exporta primitivas explícitas de control:
\begin{equation}
T_{\text{OMP}} \times T_{\text{BLAS}} \le N_{\text{physical\_cores}}
\end{equation}
controladas mediante \texttt{polydim\_set\_blas\_num\_threads(n)} y \texttt{polydim\_set\_omp\_num\_threads(n)}.

\subsection{Demostración Empírica del Límite de Roofline (Task task-940)}
Durante las pruebas asintóticas en silicio real ($D = 10{,}000, K = 512$), se demostró empíricamente la asfixia del ancho de banda DRAM en bucles sin empaquetamiento de registros:


\centering
\begin{tabular}{lrrl}
\hline
\textbf{Algoritmo} & \textbf{Tiempo Real (Silicio)} & \textbf{Error de Ortogonalidad} & \textbf{Régimen de Hardware} \\
\hline
\texttt{Cayley-SMW} (bucles planos) & $100.11$ s & $4.44 \times 10^{-16}$ & Memory-Bound (DRAM bottleneck) \\
\texttt{CholQR2} (2 pasadas Cholesky) & $9.03$ s & $8.88 \times 10^{-16}$ & $11.1\times$ Aceleración (Menos barridos) \\
\texttt{BLAS dsyrk} (V772 empaquetado) & $< 0.15$ s & $4.44 \times 10^{-16}$ & Compute-Bound (SIMD Registers) \\
\hline
\end{tabular}




Este resultado confirma que el código C++ es aritméticamente invulnerable ($\|drift\| \approx 10^{-16}$), y justifica formalmente la adopción de BLAS Nivel 3 para quebrar el muro de la memoria DRAM.
